\documentclass[12pt]{article}
\usepackage[utf8]{inputenc}
\usepackage[spanish]{babel}
\usepackage{graphicx} % Para las imágenes

% Título y Autores
\title{ \textbf{Turing y el desciframiento de Enigma: Mecanización de la lógica contra la criptografía alemana}}
\author{ \textit{\textbf{Gabriel Andrés Pla Lasa}} \\ \textit{\textbf{Carlos Daniel Largacha Leal}} \\ \\ \textit{Historia de la Computación} \\ \textit{4to Año}}
\date{\today}

\begin{document}

\maketitle % Sección 1: Portada

\begin{abstract}
    Este documento presenta una síntesis analítica e histórica sobre cómo Alan Turing revolucionó el criptoanálisis durante la Segunda Guerra Mundial, aplicando métodos de lógica matemática y modelado estadístico para vencer al sistema criptográfico de la máquina Enigma. Al explotar vulnerabilidades arquitectónicas del \textit{hardware} alemán --particularmente el hecho de que ninguna letra se cifraba sobre sí misma--, junto con fallos sistemáticos y predecibles en las rutinas de los operadores humanos, Turing y su equipo en Bletchley Park diseñaron herramientas decisivas como \textit{La Bombe} y formularon técnicas probabilísticas pioneras como el \textit{Banburismus}. Estas invenciones no solo supusieron la mecanización de la deducción lógica superando la inviabilidad de los ataques por fuerza bruta, sino que lograron reducir significativamente la duración de la guerra e influir innegablemente en la victoria aliada, al mismo tiempo que cimentaban las bases fundacionales de las Ciencias de la Computación modernas.
    
    \vspace{0.5cm}
    \textbf{Palabras clave:} Alan Turing, Enigma, La Bombe, Banburismus, Criptoanálisis, Bletchley Park.
\end{abstract}

\section{Introducción}

% Contexto de la Segunda Guerra Mundial y la supremacía de la criptografía alemana
% Introducir la figura de Alan Turing y la importancia de su trabajo en la Cabaña 8

Alan Turing fue una figura fundamental en el desciframiento de Enigma \cite{turing_enigma}.
 % Sección 2
\section{El funcionamiento de la máquina Enigma y su ``invulnerabilidad''}

% Explicar brevemente el sistema de rotores
% Aspecto Técnico Crucial: La regla de que una letra nunca se cifraba como sí misma
       % Sección 3.1
\subsection{El factor humano: Las ``Cribs'' y fallos de procedimiento}

Pese a que la restricción del reflector reducía significativamente el espacio de búsqueda, la magnitud teórica de variables seguía siendo tan abrumadora que sobrepasaba cualquier capacidad de resolución intuitiva. Incluso con esa debilidad, el sistema generaba una cantidad de combinaciones diarias que hacía inútil la prueba de ensayo y error. Sin embargo, esta robustez técnica se vio comprometida por el factor humano: la rigidez de los protocolos militares y la recurrencia de errores procedimentales por parte de los operadores alemanes introdujeron patrones predecibles que permitieron vulnerar el sistema \cite{copeland2006colossus}

Dado que las directivas militares nazis eran operadas en un marco profundamente regulado y burocrático, las comunicaciones de los artilleros, pilotos o comandantes de marina no estaban libres de redundancias dogmáticas predecibles en la redacción ordinaria del texto claro interceptado. Las rígidas doctrinas de transmisiones exigían la provisión constante e idéntica de informes meteorológicos como ``\textit{Wetterbericht}'', la enumeración regular de listados de material en franjas horarias consistentes, formulaciones invariables de localización táctica y, en forma sumaria y notoria, obligaban a las tropas a rematar religiosamente muchos comunicados solemnes con frases imperativas del partido, típicamente el ubicuo remate ``\textbf{HEIL HITLER}'' \cite{copeland2006colossus}. 

A estas predicciones o pistas basadas en el conocimiento de la jerga y rutinas militares alemanas, el equipo de Bletchley Park les asignó el término técnico \textbf{\textit{Crib}} (jerga para guía estructural de decodificación) \cite{copeland2004essential}. Bajo el enfoque computacional de Alan Turing, el \textit{Crib} operaba como un ancla lingüística que permitía ejecutar la técnica de deslizamiento o ``\textit{crib-dragging}'', el cual consistía en desplazar un fragmento de texto plano probable, como el común \textit{HEILHITLER}, sobre el flujo del criptograma \cite{copeland2006colossus}. Este proceso explotaba la vulnerabilidad crítica del diseño del reflector de Enigma: la imposibilidad física de que una letra se cifrara como ella misma ($f(x) \neq x$). Al detectar una coincidencia entre una letra del \textit{crib} y la posición correlativa en el cifrado (un ``\textit{crash}''), el analista descartaba instantáneamente esa configuración de rotores como inválida \cite{copeland2004essential}. 

Esta sinergia analítica entre los patrones predecibles del error humano y las restricciones deterministas del \textit{hardware} alemán proporcionó la lógica necesaria para alimentar el descifrado masivo y automatizado mediante el escaneo paralelo de la Bombe.        % Sección 3.2
\section{La Bombe: Mecanización del pensamiento lógico}

% El rediseño del dispositivo polaco junto a Gordon Welchman
% Análisis Técnico: Explicar la Lógica de contradicción
        % Sección 3.3
\section{El Banburismus: Estadística Bayesiana aplicada}

% Explicar cómo Turing utilizó la probabilidad para reducir el espacio de búsqueda
% El impacto en la eficiencia y el ahorro de tiempo
  % Sección 3.4
\section{Impacto Histórico y en las Ciencias de la Computación}

% Impacto bélico: El acortamiento de la guerra y el éxito en la Batalla del Atlántico
% Impacto en la CC: Cómo este trabajo validó sus teorías sobre la Máquina Universal de Turing
      % Sección 4
\section{Conclusiones}

% Resumir la relevancia de pasar de la "fuerza bruta" a la "mecanización de la lógica"
 % Sección 5

\bibliographystyle{plain}
\bibliography{referencias}          % Sección 6

\end{document}
